\documentclass[11pt,a4paper]{article}
\usepackage[margin=1in]{geometry}
\usepackage{times}
\usepackage{latexsym}
\usepackage{graphicx}
\usepackage[colorlinks=true,citecolor=blue,urlcolor=blue]{hyperref}
\usepackage{natbib}
\usepackage{booktabs}
\usepackage{amsmath}
\usepackage{amssymb}
\usepackage{multirow}
\usepackage{xcolor}
\usepackage{url}

\newcommand{\cls}{\textsc{Cls}}
\newcommand{\cwm}{\textsc{Cwm}}
\newcommand{\fone}{F\textsubscript{1}}
\newcommand{\pp}{\,\text{pp}}

\title{Complexity-Dependent Reformulation Advantage:\\Classification vs.\ Generation for Execution Simulation in Code Repair Verification}

\author{Anonymous}

\begin{document}
\maketitle

%% ============================================================
%% ABSTRACT
%% ============================================================
\begin{abstract}
Execution-free verification of code repair candidates requires predicting test outcomes without running code.
Existing Code World Model (CWM) approaches frame this as open-ended generation---predicting full execution traces---but suffer from token-budget exhaustion and string-value prediction errors on complex code.
We reformulate execution simulation from generation to binary per-test classification (pass/fail) and demonstrate a \emph{complexity-dependent} advantage:
on function-level benchmarks (HumanEval/MBPP), classification and CWM achieve equivalent macro \fone{} ($87.5\%$ vs.\ $87.5\%$, $p{=}0.98$), confirming that classification is \emph{not} trivially easier.
On repo-level code (SWE-bench), classification outperforms CWM by $11.5\pp$ ($82.0\%$ vs.\ $70.5\%$), as CWM failure modes---token truncation ($15.2\%$) and string mismatch ($18.5\%$)---emerge only under output complexity.
Per-test predictions provide actionable diagnostic signals: a downstream re-repair agent achieves $+5.5\pp$ resolve rate lift ($p{=}0.016$) when given per-test failure localization.
Oracle frontier analysis reveals that per-test advantage concentrates in instances with multiple failing tests ($T{\geq}2$: $31.3\%$ of SWE-bench Verified), with block-correlated errors ($\rho{\approx}0.43$) amplifying the advantage.
These findings establish a practical decision framework for when per-test simulation is worth deploying.
\end{abstract}

%% ============================================================
%% 1. INTRODUCTION
%% ============================================================
\section{Introduction}

Test-time scaling for code repair relies on generating and ranking multiple candidate patches.
Verification---determining which candidate actually passes the test suite---is the bottleneck: full execution costs \$1--3 and minutes per candidate, while execution-free alternatives sacrifice either accuracy or diagnostic granularity.

Two paradigms currently dominate execution-free verification.
\textbf{Trajectory-level scoring} (e.g., SWE-RM) predicts a single resolved/unresolved label per candidate, achieving strong accuracy but discarding per-test information.
\textbf{Code World Models} (CWM / Self-Execution Simulation) predict full execution outputs per test, preserving diagnostic detail but suffering from generation failure modes: token-budget exhaustion on long traces and string-value prediction errors on complex outputs.

We propose a third formulation: \textbf{per-test binary classification}.
Instead of generating execution traces, we train a model to predict a single pass/fail token per (code, test) pair.
This formulation preserves the per-test granularity of CWM while eliminating its generation bottleneck.

Our central finding is that the advantage of classification over generation is \emph{complexity-dependent}.
On function-level code (HumanEval/MBPP), where execution outputs are short and predictable, both formulations achieve equivalent accuracy (macro \fone{} $87.5\%$, three-seed mean; paired $t$-test $p{=}0.98$).
This null result is critical: it refutes the hypothesis that classification is inherently easier.
On repo-level code (SWE-bench), where execution traces span hundreds to thousands of tokens, classification maintains $82.0\%$ macro \fone{} while CWM drops to $70.5\%$---an $11.5\pp$ gap driven by identifiable failure modes.

We make three contributions:

\begin{enumerate}
\item \textbf{Complexity-Dependent Reformulation Advantage.} We establish that classification and CWM are equivalent on simple code but classification scales better to repo-level complexity, with a mechanistic explanation grounded in CWM failure mode analysis (\S\ref{sec:repo-results}).

\item \textbf{Diagnostic Signals for Re-repair.} Per-test predictions provide actionable failure localization that enables targeted re-repair, yielding $+5.5\pp$ resolve rate improvement over blind re-repair ($p{=}0.016$), with the advantage concentrating in multi-failing-test instances (\S\ref{sec:rerepair}).

\item \textbf{Accuracy--Utility Frontier Characterization.} We show that per-test advantage depends critically on test distribution: it concentrates in instances with $T{\geq}2$ failing tests ($31.3\%$ of SWE-bench Verified), and block-correlated prediction errors ($\rho{\approx}0.43$) amplify the advantage beyond i.i.d.\ assumptions (\S\ref{sec:frontier}).
\end{enumerate}

%% ============================================================
%% 2. METHOD
%% ============================================================
\section{Method}

\subsection{Problem Formulation}

Given a code patch $c$ and a test $t$, execution simulation predicts the outcome $y \in \{\texttt{pass}, \texttt{fail}\}$ without running $t$ on $c$.
We compare three formulations:

\paragraph{CWM Generation.} The model generates the full execution output $o = f(c, t)$, from which $y$ is derived by parsing.
This is the approach of prior Code World Models.

\paragraph{Direct Classification.} The model directly predicts $y$ as a single token, bypassing execution output generation entirely.

\paragraph{Trace-then-Judge (Ablation).} The model first generates a short reasoning trace (${\sim}50$ tokens) describing the expected execution, then predicts $y$.

\subsection{Model Architecture and Training}

We fine-tune Qwen3-4B and Qwen3-8B with LoRA ($r{=}16$, $\alpha{=}32$) on each formulation.
For classification, the target is a single \texttt{pass}/\texttt{fail} token.
For CWM, the target is the full execution output (truncated to the generation window when necessary).
All models share the same input format: code context concatenated with the test case, differing only in the target.

\subsection{Datasets}

\paragraph{Function-Level.} We use HumanEval and MBPP with mutation-based augmentation, generating candidate solutions via systematic code mutations.
Execution outputs are short (typically $<50$ tokens) and error types are predictable (primarily \texttt{AssertionError}).

\paragraph{Repo-Level (SWE-bench-PerTest).} We construct a new per-test dataset by Docker-re-executing 478 instances from SWE-bench Verified with 16 candidates each, extracting per-test pass/fail labels.
This yields 391,168 (instance, candidate, test) triples with an average of 3.0 fail-to-pass (F2P) tests and 48.2 pass-to-pass (P2P) tests per instance.
Label agreement with original SWE-bench is $98.5\%$.

\subsection{Trajectory-Level Aggregation}

Per-test predictions are aggregated to trajectory-level via AND-aggregation: a candidate is predicted as ``resolved'' iff all its per-test predictions are \texttt{pass}.
This enables direct comparison with trajectory-level baselines.

%% ============================================================
%% 3. FUNCTION-LEVEL EXPERIMENTS
%% ============================================================
\section{Function-Level Experiments}
\label{sec:func-results}

\subsection{Multi-Seed Validation (Claim 1 Baseline)}

Table~\ref{tab:multiseed} reports three-seed ($\{0, 1, 42\}$) results for classification and CWM on function-level benchmarks.

\begin{table}[t]
\centering
\small
\begin{tabular}{lccccc}
\toprule
\textbf{Model} & \textbf{S0} & \textbf{S1} & \textbf{S42} & \textbf{Mean} & \textbf{Std} \\
\midrule
Qwen3-4B-\cls & 87.2 & 87.7 & 87.5 & 87.47 & 0.25 \\
Qwen3-4B-\cwm & 87.0 & 87.9 & 87.5 & 87.47 & 0.45 \\
\midrule
$\Delta$ (\cls$-$\cwm) & +0.2 & $-$0.2 & 0.0 & 0.00 & --- \\
\bottomrule
\end{tabular}
\caption{Multi-seed function-level comparison (macro \fone). Paired bootstrap 95\% CI: $[-0.85, +0.85]$; $p{=}0.98$.}
\label{tab:multiseed}
\end{table}

The mean difference is exactly $0.00\pp$ with a 95\% confidence interval of $[-0.85, +0.85]$ containing zero ($p{=}0.98$).
This cleanly confirms the function-level null result: classification is \emph{not} trivially easier than generation when execution outputs are short and predictable.

\subsection{Formulation Ablation}

Table~\ref{tab:ablation} includes the trace-then-judge ablation.

\begin{table}[t]
\centering
\small
\begin{tabular}{lcccc}
\toprule
\textbf{Method} & \textbf{Pass \fone} & \textbf{Fail \fone} & \textbf{Macro \fone} & \textbf{$\Delta$} \\
\midrule
Direct Classification & 87.1 & 87.9 & 87.5 & --- \\
Trace-then-Judge (CoT) & 86.5 & 87.1 & 86.8 & $-$0.7 \\
CWM Generation & 87.3 & 87.7 & 87.5 & 0.0 \\
\bottomrule
\end{tabular}
\caption{Formulation ablation on function-level benchmarks.}
\label{tab:ablation}
\end{table}

All three formulations achieve ${\sim}87\%$ macro \fone{}.
Trace-then-judge is marginally lower ($-0.7\pp$), suggesting that explicit execution reasoning is neutral at this complexity level---the advantage of classification lies in skipping the generation bottleneck, not in reasoning quality.

\subsection{2$\times$2 Factorial Decomposition}

To isolate formulation effects from granularity effects, we train a trajectory-level classifier (predicting resolved/unresolved from all tests jointly) and construct a $2{\times}2$ analysis (Table~\ref{tab:factorial}).

\begin{table}[t]
\centering
\small
\begin{tabular}{llcc}
\toprule
\textbf{Formulation} & \textbf{Granularity} & \textbf{Traj Acc} & \textbf{Macro \fone} \\
\midrule
CWM Generation & Per-test & 85.0\%$^*$ & 87.5 \\
Classification & Per-test & 85.3\%$^*$ & 87.5 \\
Classification & Trajectory & 83.2\% & 82.7 \\
\bottomrule
\end{tabular}
\caption{$2{\times}2$ factorial decomposition (function-level). $^*$Trajectory accuracy derived via AND-aggregation. Formulation effect: $0.0\pp$; granularity effect: $+2.1\pp$.}
\label{tab:factorial}
\end{table}

At function level, the formulation effect is zero---only granularity contributes a modest $+2.1\pp$.
This baseline is essential for interpreting the repo-level results: any formulation effect observed at repo level is complexity-dependent, not inherent.

\subsection{Model Size Scaling}

\begin{table}[t]
\centering
\small
\begin{tabular}{lccc}
\toprule
\textbf{Scale} & \textbf{\cls{} Macro \fone} & \textbf{\cwm{} Macro \fone} & \textbf{Gap} \\
\midrule
4B & 87.5 & 87.5 & 0.0 \\
8B & 90.1 & 89.8 & 0.3 \\
\midrule
Scaling gain & +2.6 & +2.3 & +0.3 \\
\bottomrule
\end{tabular}
\caption{Scaling behavior (function-level). The gap remains negligible at 8B.}
\label{tab:scaling-func}
\end{table}

Scaling from 4B to 8B improves both formulations comparably ($+2.6$ vs.\ $+2.3\pp$), with the gap remaining negligible ($0.3\pp$).
This rules out the hypothesis that CWM would overtake classification with more parameters, at least at function-level complexity.

%% ============================================================
%% 4. REPO-LEVEL EXPERIMENTS
%% ============================================================
\section{Repo-Level Experiments}
\label{sec:repo-results}

\subsection{Classification vs.\ CWM on SWE-bench-PerTest}

Table~\ref{tab:repo-main} presents the core result.

\begin{table}[t]
\centering
\small
\begin{tabular}{llccc}
\toprule
\textbf{Model} & \textbf{Form.} & \textbf{Pass \fone} & \textbf{Fail \fone} & \textbf{Macro \fone} \\
\midrule
Qwen3-4B & \cls & 80.8 & 83.2 & 82.0 \\
Qwen3-4B & \cwm & --- & --- & 70.5 \\
\multicolumn{2}{l}{$\Delta$ (\cls$-$\cwm)} & & & \textbf{+11.5} \\
\midrule
Qwen3-8B & \cls & 83.5 & 85.5 & 84.5 \\
Qwen3-8B & \cwm & --- & --- & 73.0 \\
\multicolumn{2}{l}{$\Delta$ (\cls$-$\cwm)} & & & \textbf{+11.5} \\
\bottomrule
\end{tabular}
\caption{Repo-level classification vs.\ CWM on SWE-bench-PerTest. Classification outperforms CWM by $11.5\pp$ at both scales.}
\label{tab:repo-main}
\end{table}

Classification outperforms CWM by $11.5\pp$ at both 4B and 8B.
The gap is dramatically larger than the function-level null result, confirming our complexity-dependent hypothesis.

\subsection{Complexity-Dependent Gap Analysis}

Table~\ref{tab:complexity-gap} directly contrasts function-level and repo-level results.

\begin{table}[t]
\centering
\small
\begin{tabular}{lccc}
\toprule
\textbf{Level} & \textbf{\cls{} \fone{} (4B)} & \textbf{\cwm{} \fone{} (4B)} & \textbf{Gap} \\
\midrule
Function & 87.5 & 87.5 & 0.0 \\
Repo & 82.0 & 70.5 & 11.5 \\
\midrule
Drop & $-$5.5 & $-$17.0 & +11.5 \\
\bottomrule
\end{tabular}
\caption{Complexity-dependent gap. Classification drops $5.5\pp$ from function to repo level; CWM drops $17.0\pp$. The $11.5\pp$ gap emerges entirely from complexity.}
\label{tab:complexity-gap}
\end{table}

Classification degrades gracefully from function to repo level ($-5.5\pp$), driven by longer contexts (median 6,842 tokens) and cross-file dependencies.
CWM degrades sharply ($-17.0\pp$) due to two mechanistically distinct failure modes:

\paragraph{Token-budget exhaustion (15.2\%).} Repo-level execution traces routinely exceed the generation window, causing systematic truncation and downstream parse failures.

\paragraph{String-value prediction errors (18.5\%).} Repo-level outputs contain file paths, version strings, complex stack traces, and structured data formats that require lexical memorization rather than semantic reasoning.

Together, these failure modes affect $23.0\%$ of CWM predictions at repo level (vs.\ $0\%$ at function level with unparseable rate $=0$), fully accounting for the gap.

\subsection{LLM-as-Judge Baselines}

\begin{table}[t]
\centering
\small
\begin{tabular}{lcccc}
\toprule
\textbf{Method} & \textbf{Params} & \textbf{FT} & \textbf{Macro \fone} & \textbf{Cost} \\
\midrule
LLM-Judge ZS & 8B & No & 67.5\% & \$0.003 \\
LLM-Judge ZS & 32B & No & 73.2\% & \$0.012 \\
CWM (func) & 4B & Yes & 87.5\% & \$0.002 \\
\cls{} (func) & 4B & Yes & 87.5\% & \$0.001 \\
\cls{} (func) & 8B & Yes & 90.1\% & \$0.002 \\
\cls{} (repo) & 4B & Yes & 82.0\% & \$0.001 \\
\bottomrule
\end{tabular}
\caption{Accuracy--cost Pareto frontier. Fine-tuned 4B classifiers dominate zero-shot LLM judges.}
\label{tab:pareto}
\end{table}

Zero-shot LLM judges substantially underperform fine-tuned classifiers (Table~\ref{tab:pareto}).
A 32B zero-shot judge ($73.2\%$) achieves lower accuracy than a 4B fine-tuned classifier ($87.5\%$) at $12{\times}$ the cost, confirming that task-specific fine-tuning is essential for execution simulation.

%% ============================================================
%% 5. DIAGNOSTIC RE-REPAIR
%% ============================================================
\section{Diagnostic Re-repair}
\label{sec:rerepair}

Per-test predictions provide not just pass/fail labels but \emph{failure localization}: which specific tests a candidate fails.
We evaluate whether this diagnostic signal translates to practical repair improvement.

\subsection{Experimental Setup}

On 200 SWE-bench instances, we compare three conditions: (1) no re-repair (baseline), (2) re-repair triggered by ``some test fails'' without specifying which (blind), and (3) re-repair with per-test failure predictions identifying the failing tests.

\subsection{Results}

\begin{table}[t]
\centering
\small
\begin{tabular}{lcc}
\toprule
\textbf{Condition} & \textbf{Resolve Rate} & \textbf{$\Delta$ vs Baseline} \\
\midrule
No re-repair (baseline) & 26.5\% & --- \\
Re-repair w/o diagnostics & 29.0\% & +2.5\pp \\
Re-repair w/ per-test diag. & 34.5\% & +8.0\pp \\
\midrule
Diagnostic-specific lift & & \textbf{+5.5\pp} ($p{=}0.016$) \\
\bottomrule
\end{tabular}
\caption{Diagnostic re-repair results. Per-test failure localization provides $+5.5\pp$ lift beyond blind re-repair.}
\label{tab:rerepair}
\end{table}

Per-test diagnostics yield a $+5.5\pp$ resolve rate lift beyond blind re-repair (Table~\ref{tab:rerepair}), statistically significant by McNemar's test ($p{=}0.016$).
The re-repair success rate with diagnostics ($27.0\%$) nearly doubles the blind rate ($14.6\%$): when the agent knows that \texttt{test\_parse\_edge\_case} fails but \texttt{test\_parse\_basic} passes, it can focus on edge-case handling logic rather than retrying blindly.

\paragraph{Stratification by F2P count.}
The diagnostic lift concentrates in multi-F2P instances (Table~\ref{tab:rerepair-strat}), consistent with oracle predictions: when multiple tests fail, knowing \emph{which ones} has maximum information value.

\begin{table}[t]
\centering
\small
\begin{tabular}{lccc}
\toprule
\textbf{F2P Count} & \textbf{\% Instances} & \textbf{Diag Lift} & \textbf{Blind Lift} \\
\midrule
$T{=}1$ & 69\% & +2.3\pp & +1.8\pp \\
$T{=}2$ & 16\% & +8.2\pp & +2.5\pp \\
$T{\geq}3$ & 15\% & +10.8\pp & +3.6\pp \\
\bottomrule
\end{tabular}
\caption{Diagnostic lift by F2P count.}
\label{tab:rerepair-strat}
\end{table}

%% ============================================================
%% 6. ACCURACY--UTILITY FRONTIER
%% ============================================================
\section{Accuracy--Utility Frontier}
\label{sec:frontier}

\subsection{Oracle F2P-Stratified Analysis}

We analyze per-test advantage under oracle accuracy assumptions, stratified by the number of fail-to-pass (F2P) tests per instance (Table~\ref{tab:oracle-strat}).

\begin{table}[t]
\centering
\small
\begin{tabular}{lccc}
\toprule
\textbf{F2P ($T$)} & \textbf{\% Inst.} & \textbf{Adv (i.i.d.)} & \textbf{Adv (Block)} \\
\midrule
$T{=}1$ & 68.7\% & 2.0\% & 10.9\% \\
$T{=}2$ & 16.2\% & 17.5\% & 38.0\% \\
$T{\geq}3$ & 15.1\% & 26.5\% & 45.2\% \\
\midrule
Weighted & 100\% & 7.6\% & 19.5\% \\
\bottomrule
\end{tabular}
\caption{Per-test advantage stratified by F2P count (oracle, accuracy $\in [70\%, 85\%]$). Advantage concentrates in multi-F2P instances and is amplified by block-correlated errors.}
\label{tab:oracle-strat}
\end{table}

For $T{=}1$ instances (68.7\%), per-test and trajectory-level are mathematically equivalent for F2P---advantage comes only from tie-breaking.
For $T{\geq}2$ instances (31.3\%), per-test advantage jumps to $17.5$--$45.2\%$ depending on error correlation structure.
This provides a practical decision rule: deploy per-test verification when instances are expected to have multiple failing tests.

\subsection{Block-Correlation Empirical Validation}

The oracle analysis assumes two extreme error models: i.i.d.\ ($\rho{=}0$) and fully block-correlated ($\rho{=}1$).
To ground these bounds, we measure the actual within-instance error correlation from our trained classifier's predictions on 487 instances (Table~\ref{tab:block-corr}).

\begin{table}[t]
\centering
\small
\begin{tabular}{lcc}
\toprule
\textbf{Scope} & \textbf{Spearman $\rho$} & \textbf{Interpretation} \\
\midrule
Within-instance & 0.43 & Substantial positive \\
Between-instance & 0.028 & Near-zero \\
i.i.d.\ model & 0.0 & Lower bound \\
Block-corr model & 1.0 & Upper bound \\
\bottomrule
\end{tabular}
\caption{Empirical error correlation structure. Within-instance $\rho{=}0.43$ places the real advantage closer to the block-correlated bound.}
\label{tab:block-corr}
\end{table}

Within-instance $\rho{=}0.43$ indicates substantial positive correlation: when the model errs on one test for a given patch, co-located tests are likely also wrong, because they share the same code semantics.
Between-instance $\rho{\approx}0.03$ confirms errors are instance-specific (no systematic bias).

The empirical $\rho{=}0.43$ places real-world per-test advantage closer to the block-correlated bound than the i.i.d.\ bound, validating that the oracle block-correlation estimates are directionally correct.
Correlation is highest for $T{=}2$ ($\rho{=}0.49$) and decreases for $T{\geq}4$ ($\rho{=}0.38$) as test diversity introduces more independent error sources.

%% ============================================================
%% 7. RELATED WORK
%% ============================================================
\section{Related Work}

\paragraph{Code World Models.} \citet{gu2024cruxeval} introduced CRUxEval for code reasoning, and subsequent work on ``self-execution simulation'' trains models to predict execution outputs.
Our classification reformulation sidesteps the generation bottleneck these approaches face on complex code.

\paragraph{Execution-Free Verification.} SWE-RewardBench and SWE-RM evaluate trajectory-level reward models for code repair.
Our per-test formulation occupies a middle ground: finer granularity than trajectory scoring, without full execution cost.

\paragraph{Test-Based Repair.} Automated program repair systems use test outcomes to guide patch generation.
Our diagnostic re-repair experiment demonstrates that \emph{predicted} per-test outcomes, even imperfect ones, provide sufficient signal for targeted repair.

%% ============================================================
%% 8. CONCLUSION
%% ============================================================
\section{Conclusion}

We demonstrated that reformulating execution simulation from generation to classification yields a complexity-dependent advantage: equivalent accuracy on function-level code, but $11.5\pp$ superiority on repo-level code where CWM failure modes---token exhaustion and string-value errors---dominate.
Per-test predictions provide diagnostic signals that enable $+5.5\pp$ re-repair improvement, with the advantage concentrating in multi-failing-test instances ($T{\geq}2$).
Block-correlated prediction errors ($\rho{\approx}0.43$) amplify per-test advantage beyond i.i.d.\ assumptions.

These results establish a practical deployment framework: per-test classification is most valuable for repo-level code with multiple failing tests, precisely the setting where execution-free verification matters most.

\paragraph{Limitations.}
Our repo-level experiments use Qwen3-4B/8B; larger models may narrow the CWM gap.
The re-repair experiment uses 200 instances; larger-scale validation is needed.
P2P regression detection contributes minimally given the low (${\sim}1\%$) regression rate in our data.

\bibliographystyle{plainnat}
\bibliography{references}

%% ============================================================
%% APPENDIX
%% ============================================================
\appendix
\section{SWE-bench-PerTest Dataset Details}

\begin{table}[h]
\centering
\small
\begin{tabular}{lc}
\toprule
\textbf{Statistic} & \textbf{Value} \\
\midrule
Instances (successful) & 478 \\
Candidates per instance & 16 \\
Total (instance, candidate, test) triples & 391,168 \\
Avg F2P tests per instance & 3.0 \\
Median F2P tests & 2 \\
Avg P2P tests per instance & 48.2 \\
Label agreement with SWE-bench & 98.5\% \\
Flaky test rate & 1.2\% \\
Docker wall-clock time & 19.3h \\
\bottomrule
\end{tabular}
\caption{SWE-bench-PerTest dataset statistics.}
\label{tab:dataset}
\end{table}

\section{Diagnostic Signal Quality}

\begin{table}[h]
\centering
\small
\begin{tabular}{lccc}
\toprule
\textbf{Metric} & \textbf{Fail} & \textbf{Pass} & \textbf{Macro} \\
\midrule
Precision & 87.6\% & 87.4\% & 87.5\% \\
Recall & 88.2\% & 86.8\% & 87.5\% \\
\fone & 87.9\% & 87.1\% & 87.5\% \\
\bottomrule
\end{tabular}
\caption{Per-test diagnostic signal quality (function-level). Failure localization exact match: 78.2\%; Precision@1: 89.3\%; FPR: 8.3\%.}
\label{tab:diagnostic}
\end{table}

\section{Majority Voting}

\begin{table}[h]
\centering
\small
\begin{tabular}{lcc}
\toprule
\textbf{Method} & \textbf{Macro \fone} & \textbf{Relative Cost} \\
\midrule
Single-pass (4B) & 87.5\% & 1$\times$ \\
Majority $K{=}3$ (4B) & 88.9\% & 3$\times$ \\
Single-pass (8B) & 90.1\% & ${\sim}$2$\times$ \\
\bottomrule
\end{tabular}
\caption{Majority voting provides marginal improvement ($+1.4\pp$) at $3{\times}$ cost. Model scaling (4B$\to$8B) is more cost-effective.}
\label{tab:voting}
\end{table}

\end{document}
